\section{Conventions and preliminary facts}\label{sec:conventions}

\subsection{Measures, scales and directions}

All measures are finite Borel measures. A subprobability has mass at most
one. Restrictions are left unnormalized unless normalization is explicitly
specified. We write
\[
 I_s(\mu)=\iint |x-y|^{-s}\,\dd\mu(x)\dd\mu(y),
 \qquad
 \widehat\mu(\xi)=\int e^{-i x\cdot\xi}\,\dd\mu(x).
\]
For measures on the line we use the same Fourier convention.
The norm \(\|\cdot\|_{\TV}\) is total variation.
The notation \(f\lesssim g\) means \(f\le Cg\), with the dependence of \(C\)
specified by the ambient discussion.

A measure has an \(s\)-ball bound with loss \(K\), down to scale
\(\delta\), if
\[
 \mu(B(z,r))\le K r^s
 \quad\text{for every }z\text{ and }\delta\le r\le1.
\]
A cap bound on a sphere has the analogous meaning for its Euclidean
chordal metric. Passing between this metric, angular distance and fixed
smooth charts changes constants only. We may use dyadic radii, and enlarge
the constant to recover all radii in the specified interval.

For \(x\ne y\), put
\[
 \pi_x(y)=\frac{y-x}{|y-x|},\qquad D(x,y)=|x-y|.
\]
A tube is a neighborhood of an affine line; a plate is a neighborhood
of an affine plane of the dimension indicated. Directions are oriented
unless an antipodal pair of caps is explicitly used.
All configurations lie in bounded sets. Fixed rescalings and a finite
number of coordinate charts have constant cost.

The original ambient dimension is always denoted by \(d\), and
\(S=d/2\). We use \(R=2^N\) for the reciprocal final scale.
In odd dimension we write \(d=2k-1\); the linear angular gain is then used
in \(\R^{k+1}\), while \(S=k-1/2\) still refers to the original dimension.
Letters denoting dimensions inside a preliminary lemma are local to that
lemma.

\subsection{Finite energy and projection}

We record elementary consequences of the usual Frostman and projection
principles, including the restriction operation needed later.

\begin{lemma}\label{conv:energy}
Let \(\mu\) be supported in a bounded subset of \(\R^d\).
\begin{enumerate}[label=(\roman*)]
\item If \(\mu(B(z,r))\le C r^s\) for \(0<r\le1\), then
\(I_t(\mu)<\infty\) for every \(0<t<s\).
\item If \(I_t(\mu)<\infty\), then for every \(\eps>0\) there is a
compact restriction of mass at least \(\mu(\R^d)-\eps\) that satisfies
a \(t\)-ball bound with a finite constant.
\item If \(0<t<m\le d\), \(I_t(\mu)<\infty\), and \(P_V\) is
orthogonal projection to a Haar-random \(m\)-plane, then
\(I_t((P_V)_\#\mu)<\infty\) for almost every \(V\).
\end{enumerate}
\end{lemma}

\begin{proof}
For (i), split the integral at distance one and sum dyadic annuli.
The contribution at distance at most \(2^{-j}\) is bounded by
\(C2^{-j(s-t)}\), up to a constant depending on the total mass.

For (ii), let \(U_t^\mu(x)=\int |x-y|^{-t}\,\dd\mu(y)\).
This is finite for \(\mu\)-almost every \(x\), so a sufficiently large
level \(M\) discards arbitrarily little mass. Choose a compact subset
\(K\subset\{U_t^\mu\le M\}\) losing at most the remaining prescribed mass.
If \(K\cap B(z,r)\ne\varnothing\), choose \(x\) in this intersection.
Then
\[
 \mu|_K(B(z,r))
 \le\mu(B(x,2r))\le (2r)^t U_t^\mu(x)\le M(2r)^t.
\]
The same bound is immediate if the intersection is empty.

For (iii), rotational invariance and the integrability of the negative
\(t\)-moment of the length of an \(m\)-dimensional projection give
\[
 \int_{G(d,m)}|P_V z|^{-t}\,\dd V\le C_{d,m,t}|z|^{-t}
 \quad(z\ne0).
\]
For \(m<d\), this follows by integrating the density of the first \(m\)
coordinates of a uniform unit vector near zero: the radial integral there
is bounded by a constant times
\(\int_0^1 r^{m-1-t}\,\dd r<\infty\). For \(m=d\), the identity projection
makes the assertion immediate. Integrating in \((x,y)\) and applying
Tonelli proves the claim.
\end{proof}

We use Frostman's Lemma in its standard form
\cite[Theorem~8.8]{Mattila1995}: if a compact set \(E\)
has \(\dim_H E>s\), it supports a nonzero finite measure satisfying an
\(s\)-ball bound at all scales. Dividing by its mass gives a probability.
The positive dimension margin permits a slight decrease of \(s\) whenever
a finite energy rather than a ball bound is needed.

Dyadic cubes are half-open, so they form genuine partitions.
One may choose a translate of the dyadic grid whose boundaries have zero
mass for the finitely many starting measures: for each coordinate
hyperplane this follows by integrating its mass over the translation,
and the set of grid hyperplanes is countable.

\subsection{Subpower losses and fixed gains}

For a sequence \(R\to\infty\), a positive factor is at most subpower if it
is bounded by \(R^{o(1)}\). A retained mass is at least subpower if it is
bounded below by \(R^{-o(1)}\). These are statements about the sequence;
their use does not assert that an arbitrary fixed positive exponent is
zero. Constants and powers of \(\log R\) are subpower.
The corresponding notation for \(\delta\to0\) uses \(R=\delta^{-1}\).

We will use the following quantifier principle. It makes precise the
compactness passages in which the input losses tend to zero and an
argument produces a fixed positive improvement.

\begin{lemma}[Sequential uniformity principle]\label{conv:uniformity}
For fixed background parameters, let \(\mathcal C(\eps,R)\) be classes
of data, increasing in the allowed loss \(\eps>0\).
Let \(\mathcal G(c,R)\) be conclusions, weakening as the gain
\(c>0\) decreases. Suppose that for every sequence
\[
 \eps_j\downarrow0,\qquad R_j\to\infty,\qquad
 X_j\in\mathcal C(\eps_j,R_j),
\]
there are a subsequence and a fixed \(c>0\) on which
\(X_j\) satisfies \(\mathcal G(c,R_j)\) for all sufficiently large \(j\).
Then there exist \(\eps_0,c_0>0\) and \(R_0<\infty\) such that
every \(X\in\mathcal C(\eps_0,R)\), \(R\ge R_0\), satisfies
\(\mathcal G(c_0,R)\).
\end{lemma}

\begin{proof}
If the conclusion fails, for each \(j\) choose \(R_j\ge j\) and
\(X_j\in\mathcal C(1/j,R_j)\) failing \(\mathcal G(1/j,R_j)\).
The assumed subsequence satisfies \(\mathcal G(c,R_j)\) for some fixed
\(c>0\). For large \(j\), \(1/j<c\), so monotonicity implies
\(\mathcal G(1/j,R_j)\), a contradiction.
\end{proof}

Fixed constants in a power-small output may be absorbed by decreasing the
output exponent and increasing the scale threshold. When several losses
occur in an input, their maximum serves as \(\eps\) in
Lemma~\ref{conv:uniformity}. Compactness of uniformly Lipschitz normalized
profiles is the usual Arzel\`a--Ascoli compactness
\cite[Theorem~4.43]{Folland1999}.
All background parameters, including any strict dimension or
transversality exponent, are fixed before this principle is applied.
The detailed arguments below specify the required ordering of losses.

For measurable selection at a fixed scale it suffices to work with
dyadic grids, finite tube nets and a finite number of charts.
Strict inequalities can first be imposed with a fixed enlargement of
widths. Thus unions of exceptional configurations can be chosen Borel;
no selection from an uncountable unparameterized family is needed.
